\documentclass[conference]{IEEEtran}
\IEEEoverridecommandlockouts
\usepackage{amsmath,amssymb}   % amssymb also provides \checkmark and loads amsfonts
\usepackage{graphicx}
\usepackage{textcomp}
\usepackage{xcolor}
\usepackage{booktabs}
\usepackage{enumitem}
\usepackage{float}             % needed for [H] placement
\usepackage{listings}
\usepackage{csquotes}          % recommended by biblatex

\lstset{
  basicstyle=\ttfamily\footnotesize,
  breaklines=true,
  columns=fullflexible,
  frame=single
}

\setlist{nosep,leftmargin=*}

% Status tags
\newcommand{\done}{\textcolor{green!40!black}{\textbf{[$\checkmark$ Done]}}}
\newcommand{\wip}{\textcolor{orange!70!black}{\textbf{[$\circ$ In progress]}}}
\newcommand{\planned}{\textcolor{blue!60!black}{\textbf{[$\square$ Planned]}}}
% Red TODO notes; remove every \todo before submission (search for "\todo").
\newcommand{\todo}[1]{\par\noindent\textcolor{red}{\textbf{TODO:} #1}\par}

\usepackage[style=ieee, backend=biber]{biblatex}
\addbibresource{../references.bib}

\title{Topic Modeling of Autonomous Vehicle News Coverage Across Publications, 2016--2020\\
\large Progress Report 2
}

\author{\IEEEauthorblockN{Lauren Gamber }
\IEEEauthorblockA{\textit{Department of Mathematical Sciences} \\
\textit{Montana State University}\\
Bozeman, MT \\
Lauren.Gamber@montana.edu}
\and
\IEEEauthorblockN{Laura Banaszewski}
\IEEEauthorblockA{\textit{Gianforte School of Computing} \\
\textit{Montana State University}\\
Bozeman, MT \\
laura.banaszewski@montana.edu}
\and
\IEEEauthorblockN{Mary Kate Carara}
\IEEEauthorblockA{\textit{Department of Mathematical Sciences} \\
\textit{Montana State University}\\
Bozeman, MT \\
connellymarykate@gmail.com}
}

\date{Due: November 8th, 2026}

\begin{document}

\maketitle

\newpage

\begin{center}\large\textbf{Implementation Progress and Initial Findings}\end{center}

Since Progress Report 1, we have (i) finalized the corpus by dropping five duplicate article bodies, keeping the earliest-dated copy of each (7,337 to 7,332 articles); (ii) completed the Introduction and Related Work sections; (iii) written shared code for corpus loading, evaluation tokenization, coherence metrics, and the joint comparison; (iv) implemented LDA preprocessing, including a 50-term custom stopword list and the document-term matrix, and fit LDA for seven topic counts; and (v) implemented the full BERTopic pipeline, fit a six-configuration grid, selected a configuration, and checked its seed stability. The BERTopic runs below were fit on all 7,337 articles before duplicate removal and will be rerun on the deduplicated corpus before Progress Report 3.

\textbf{BERTopic configuration study (preliminary).} We fit six BERTopic configurations, crossing UMAP \texttt{n\_neighbors} $\in\{15,30\}$ with HDBSCAN \texttt{min\_cluster\_size} $\in\{25,50,100\}$, all with seed 42 (Table~\ref{tab:bertopic-grid}). Smaller minimum cluster sizes produced more topics (70--71 vs.\ 24) with higher mean $c_v$ coherence and fewer unassigned articles. Cosine silhouette was low for every configuration (0.095--0.130), indicating weak separation between topics in the full embedding space.

\begin{table}[H]
\caption{BERTopic configuration study (seed 42; 7,337 articles before duplicate removal). $c_v$ is the mean over topics, using the top ten unigrams per topic. Unassigned is HDBSCAN's outlier class ($-1$). Silhouette uses cosine distance on the embeddings, excluding unassigned articles. Bold: selected configuration. $^\dagger$Shortlisted for selection review (top three by $c_v$). Preliminary.}
\label{tab:bertopic-grid}
\centering
\footnotesize
\setlength{\tabcolsep}{3pt}
\begin{tabular}{rrrrrr}
\toprule
\textbf{Neighbors} & \textbf{Min.\ cluster} & \textbf{Topics} & \textbf{$c_v$} & \textbf{Unassigned} & \textbf{Silhouette} \\
\midrule
\textbf{15}$^\dagger$ & \textbf{25} & \textbf{70} & \textbf{0.701} & \textbf{17.1\%} & \textbf{0.095} \\
15$^\dagger$ & 50  & 39 & 0.661 & 19.2\% & 0.126 \\
15           & 100 & 24 & 0.635 & 21.0\% & 0.111 \\
30$^\dagger$ & 25  & 71 & 0.706 & 18.5\% & 0.121 \\
30           & 50  & 34 & 0.652 & 18.0\% & 0.120 \\
30           & 100 & 24 & 0.632 & 21.8\% & 0.130 \\
\bottomrule
\end{tabular}
\end{table}

\textbf{Selected BERTopic configuration (preliminary).} Of the three seed-42 runs with the highest $c_v$, the model-selection step chose UMAP \texttt{n\_neighbors}$=15$ with HDBSCAN \texttt{min\_cluster\_size}$=25$ (run \texttt{n15\_c25\_s42}), which had the highest mean rating of its top five topics (3.8 on a 1--5 scale, vs.\ 3.7 for the other two). It produced 70 topics with mean $c_v=0.701$ and left 17.1\% of articles (1,256) unassigned. The \texttt{n30\_c25} run had slightly higher coherence ($c_v=0.706$) but a higher unassigned share (18.5\%). Topic labels have not yet been assigned.

To check stability, we refit the selected configuration with seeds 7 and 21 and compared each article's topic assignment with the seed-42 run (Table~\ref{tab:stability}). The adjusted Rand index (ARI) is computed over articles assigned to a topic in both runs; outlier agreement is the share of all articles that both runs either assign or leave unassigned. Both ARI values exceed 0.95, so the clustering of assigned articles is largely unchanged across seeds, although 21--23\% of articles are unassigned in at least one of the two runs being compared.

\begin{table}[H]
\caption{Seed stability of the selected BERTopic configuration (\texttt{n15\_c25}), compared with the seed-42 run. Preliminary.}
\label{tab:stability}
\centering
\footnotesize
\begin{tabular}{lrrr}
\toprule
\textbf{Seed} & \textbf{ARI} & \textbf{Articles compared} & \textbf{Outlier agreement} \\
\midrule
7  & 0.957 & 77.3\% & 91.0\% \\
21 & 0.975 & 79.2\% & 92.2\% \\
\bottomrule
\end{tabular}
\end{table}

\textbf{Selected LDA configuration (preliminary).}
We built a custom stopword list from the 100 tokens with the highest document frequency (Table~\ref{tab:stopwords}).

During $K$ selection, we initially set \texttt{max\_iter}$=20$, where \texttt{max\_iter} is the maximum number of passes LDA makes over all the articles while fitting. From these runs we initially selected $K=40$ (Table~\ref{tab:lda-grid}), but a stability check with seeds 7 and 21 showed that LDA dominant-topic assignments agreed only moderately with the seed-42 run.

\todo{[Laura] Report runtime (embedding cold vs.\ cached, fitting) and implementation challenges with attempted solutions, e.g., the repository refactor and the shared corpus manifest.}

\begin{center}\large\textbf{Paper Outline 2}\end{center}
\noindent\emph{Status legend:} \done~completed;
\wip~started, partially complete; \planned~not yet started.

\section*{Abstract}
\textbf{Purpose:} The abstract will give an overview of the themes identified in news coverage of autonomous vehicles from 2016 to 2020, and briefly compare classical and modern topic modeling pipelines.

\textbf{Planned content:} 
\begin{itemize}
    \item Problem and motivation \done
    \item Data and methods \wip
    \item Thematic findings \wip
    \item Pipeline comparison \planned
\end{itemize}

\textbf{Working draft:}
Autonomous vehicle news coverage grew rapidly between 2016 and 2020, but its dominant
themes remain poorly characterized. We apply a classical pipeline and a modern embedding-based pipeline to 7,332 articles from the ``All the News'' dataset~\cite{mckinley2021}, available on Kaggle, comparing topic coherence and interpretability. We identified...

\todo{[All] Name the methods (LDA vs.\ BERTopic with GTE-small, UMAP, HDBSCAN) and the evaluation ($c_v$ coherence, two-reviewer human evaluation, Jaccard topic overlap). Replace ``We identified...'' with clearly labeled preliminary findings and the anticipated contribution, without unsupported claims.}

\section*{Introduction}
\textbf{Purpose:} The introduction will familiarize the reader with the problem of poorly characterized themes in autonomous vehicle news coverage between 2016 and 2020, and summarize the gaps in prior research. The overall research takeaways from this paper will be summarized.

\textbf{Working draft:}
In recent years, autonomous vehicles (AVs) have received extensive news coverage spanning technological development, safety, regulation, and commercial usage. Because news coverage shapes how the public perceives emerging technologies~\cite{christen2007media, anania2018effects}, understanding what this coverage emphasizes matters for researchers, policymakers, and the public. However, AV coverage spans many news outlets and consists of large volumes of unstructured text, making recurring themes difficult to identify. Topic modeling is a method that can identify those themes. 

Recent work has applied topic modeling to AV news~\cite{kang2025unveiling}, but two questions remain open. Prior analyses pooled articles across outlets, so it
is unknown whether publications differ in which themes they emphasize. In addition, no study has applied classical and embedding-based topic models to the same AV news corpus, so it is unclear whether the reported themes depend on the
choice of model. We address these gaps with three research questions:

\begin{itemize}
    \item What themes characterize news coverage of autonomous vehicles from 2016 to 2020?
    \item How does the prevalence of these themes differ across publications?
    \item How does the choice of topic model affect the themes identified?
\end{itemize}

We apply LDA and BERTopic to 7,332 AV-related articles from 25 publications in the All the News dataset~\cite{mckinley2021}. We contribute (i) the first topic-model-based comparison of AV coverage across individual U.S. publications and (ii) a matched evaluation of LDA and BERTopic on the same corpus, testing how the choice of model changes the themes identified.


\section*{Related Work}

\textbf{Purpose:} Position our study relative to prior work on autonomous vehicle media coverage and topic modeling, and show what distinguishes it.

\textbf{Working draft:}
News media are a major source of information about emerging technologies and can shape public perceptions \cite{christen2007media}. For autonomous vehicles (AVs), this influence appears to start with the headline: Anania et al.~\cite{anania2018effects} found that participants who read positive article titles reported greater willingness to ride in a driverless vehicle than those who read negative ones.

Penmetsa et al.~\cite{penmetsa2023news} measured sentiment in U.S. AV news coverage but did not identify themes. Jelinski et al.~\cite{jelinski2021undifferentiated} examined themes, but only in German coverage. Neither study applies topic modeling.

Bakalos et al.~\cite{bakalos2020} studied public opinion of autonomous transportation using sentiment analysis of social media posts. This work shows how the public feels about AVs but does not identify which themes news coverage emphasizes.

Topic modeling can identify those themes. Classical models such as Latent Dirichlet Allocation (LDA)~\cite{blei2003lda} represent each article as a mixture of topics and each topic as a word distribution. LDA infers topics from how words co-occur in articles. LDA assumes a bag-of-words model, meaning it forms topics without considering word order. Non-negative matrix factorization (NMF)~\cite{lee1999nmf} is another classical topic modeling technique that factors a document-term matrix. These classical models are very interpretable and inexpensive to run. However, they ignore context, and the number of topics has to be chosen in advance. 

Embedding-based models instead represent articles using pretrained language models. BERTopic~\cite{grootendorst2022bertopic} is one of the most widely known embedding-based models. It embeds articles, reduces dimensionality with UMAP~\cite{mcinnes2018umap}, clusters articles with HDBSCAN~\cite{campello2013hdbscan}, and describes each cluster with class-based TF-IDF~\cite{alammar2024}. An advantage of this type of topic modeling is that the embeddings capture context, so BERTopic can group articles with similar meaning even when they use different words. Additionally, the number of topics comes from clustering. However, BERTopic assigns each article to only one topic and is more computationally expensive. 

The study closest to ours is Kang et al.~\cite{kang2025unveiling}, who applied BERTopic to 11,708 AV-related news articles published between 2002 and 2021 in the ProQuest ABI/INFORM database. By analyzing paragraphs pooled across outlets, they identified 45 AV topics, showing that the AV news coverage extends beyond safety and performance to battery supply chains, regulatory uncertainty, and infrastructure. 

Kang et al. pooled paragraphs from all sampled outlets, so their results do not investigate whether different publications emphasize different themes. Additionally, no study has fit classical and embedding-based topic models to the same AV news corpus. Because classical and embedding-based topic models differ in interpretability, stability, and cost, it is unclear whether the identified themes reflect model choice or the true topics of AV news coverage. 

Our study addresses both gaps. We compare autonomous vehicle news themes across 25 publications and apply an LDA pipeline and a BERTopic pipeline to the same corpus of news articles.


\section*{Data}
\textbf{Purpose:} Identify and familiarize the reader with the dataset, including its source, structure, and audit.

\textbf{Planned content:} 
\begin{itemize}
    \item Identify data \done
    \item Complete data audit \done
    \item Perform data filtering \wip
    \item Data statistics/figures \wip
\end{itemize}

\textbf{Working draft:}
Our research utilizes the ``All the News'' dataset \cite{mckinley2021}, available freely on Kaggle. We downloaded the dataset in September of 2026. The full dataset includes 2.79 million English news articles from January 2016 to April 2020. We used a regular expression to identify a subset of 7,337 articles. The subset covers 25 publications from January 1, 2016, through March 26, 2020. Each observation represents one article. The articles are copyrighted news text used for research only, and the only personal data is author bylines.

% Body
\begin{figure}[H]
\centering
\begin{lstlisting}
self[- ]driving|driverless|autonomous (car|vehicle|truck|driving)|robo[- ]?taxi|waymo|tesla'?s? autopilot
\end{lstlisting}
\caption{Case-insensitive regular expression used to identify AV-related articles.}
\label{fig:regex}
\end{figure}

The CSV contains 15 columns: article ID, an additional source index, date, year, month, day, author, title, article body, URL, section, publication, title match, body hits, and word count. Article bodies will be the modeling input; publication and date will be metadata for interpreting and comparing results. There is no supervised target variable, so the data does not have labels or class distribution; topics are discovered rather than predicted.

An initial audit identified missing author and section values and five duplicate article bodies, with no missing article bodies or publications. Both pipelines use the same documented article set and duplicate-handling policy: for articles with duplicate bodies, the earliest-dated article is retained and the others are dropped, bringing the article total to 7,332. Since author and section are not used for our analysis, we retained articles missing this data.


\begin{table}[H]
\caption{Corpus summary by publication after duplicate removal (N = 7,332 articles). Word counts are article medians.}
\label{tab:corpus}
\centering
\footnotesize
\setlength{\tabcolsep}{3pt}
\begin{tabular}{lrrcr}
\toprule
\textbf{Publication} & \textbf{Articles} & \textbf{\%} & \textbf{First--Last} & \textbf{Med. words} \\
\midrule
TechCrunch          & 1,214 & 16.6 & 2016-01--2020-03 & 433 \\
CNBC                &   977 & 13.3 & 2016-01--2020-03 & 456 \\
Reuters             &   854 & 11.6 & 2016-01--2020-03 & 419 \\
The Verge           &   830 & 11.3 & 2016-01--2020-03 & 505 \\
Mashable            &   554 &  7.6 & 2016-01--2020-03 & 403 \\
The New York Times  &   501 &  6.8 & 2016-01--2020-03 & 1,014 \\
Vox                 &   500 &  6.8 & 2016-01--2020-02 & 678 \\
Axios               &   360 &  4.9 & 2017-01--2019-12 & 326 \\
The Hill            &   353 &  4.8 & 2016-01--2020-02 & 737 \\
Wired               &   249 &  3.4 & 2016-01--2020-03 & 614 \\
Gizmodo             &   215 &  2.9 & 2016-01--2019-09 & 496 \\
Business Insider    &   200 &  2.7 & 2016-02--2020-03 & 653 \\
Other (13 pubs.)    &   525 &  7.2 & 2016-01--2020-03 & 691 \\
\midrule
\textbf{Total}      & \textbf{7,332} & \textbf{100} & 2016-01--2020-03 & 496 \\
\bottomrule
\end{tabular}
\end{table}

The regular expression is shown in Fig.~\ref{fig:regex}. Table~\ref{tab:corpus}
shows uneven coverage: the four largest publications contribute 52.9\% of
articles, while 13 publications together contribute 7.2\%. Article length is
right-skewed (median 496, mean 761, maximum 28{,}048 words), which motivates
chunked embeddings.

For feature construction, the classical pipeline (LDA) uses unigram bag-of-words counts, whereas BERTopic uses lightly cleaned full text for embeddings and unigram counts for keywords.

\todo{[All] Complete the 30-article inspection in \texttt{data/corpus\_review.csv} (relevance, boilerplate, malformed) and report the results.}
We will inspect a random sample for relevance, boilerplate, and malformed text.

\section*{Methodology}

\textbf{Purpose:} Describe how the topic modeling analysis of self-driving car news coverage is being performed.

\textbf{Planned content:}
\begin{itemize}
    \item Corpus construction, data audit, and duplicate handling \done
    \item Classical pipeline: LDA baseline \cite{blei2003lda} \wip
    \item Modern pipeline: BERTopic \wip
    \item Topic labeling and publication comparisons \planned
\end{itemize}

\textbf{Working draft:}
\subsection*{Overview}
We compare a classical Latent Dirichlet Allocation (LDA) baseline with a modern BERTopic pipeline. Both pipelines are applied to the same set of 7,332 articles from 25 publications (January 2016 to March 2020), after applying the shared duplicate-handling policy and preprocessing suited to each representation.

\subsection*{Classical Pipeline (LDA)}
The classical pipeline follows four stages \cite{qamar2024}: (1) data preprocessing, (2) document-term representation, (3) topic model fitting, and (4) topic interpretation. We tokenize article bodies, remove English and custom stopwords, and build a document-term count matrix. LDA provides an interpretable baseline through its topic-word distributions and article-level topic mixtures. We compare a small set of topic counts and assign each article its highest-weight (dominant) topic for publication summaries, while retaining the full topic mixture.

Tokenization was performed by (1) merging words with apostrophes and hyphens, (2) deleting non-ASCII characters, (3) dropping words with fewer than three characters, and (4) lowercasing all words.

After removing scikit-learn's English stopwords, number-only tokens, and tokens shorter than three characters, we ranked the remaining tokens by document frequency and inspected the top 100. We removed 40 reporting and filler terms (e.g., \textit{said}, \textit{told}, \textit{according}, \textit{thats}) and 10 terms that appear in 53--88\% of articles because they define the corpus (e.g., \textit{selfdriving}, \textit{autonomous}, \textit{car}, \textit{company}); the full list is in Table~\ref{tab:stopwords}. Custom stopword removal reduced the median article length from 264 to 227.5 tokens (mean 391.5 to 339.6; range 70--11,986 to 56--10,353).

\begin{table}[H]
\caption{Custom stopwords removed for the LDA pipeline, chosen from the 100 tokens with the highest document frequency.}
\label{tab:stopwords}
\centering
\footnotesize
\begin{tabular}{@{}p{0.22\columnwidth}p{0.72\columnwidth}@{}}
\toprule
\textbf{Group} & \textbf{Terms} \\
\midrule
Reporting and filler (40) & said, says, say, told, according, like, just, thats, dont, did, going, think, lot, right, way, make, making, use, used, using, need, want, come, able, called, including, big, long, far, better, end, number, new, time, year, years, week, today, early, earlier \\
\addlinespace
Corpus-defining or generic (10) & selfdriving, autonomous, cars, car, vehicles, vehicle, company, companies, companys, technology \\
\bottomrule
\end{tabular}
\end{table}

Once tokenization was complete, we represented each article as unigram counts (not TF-IDF), using scikit-learn's \texttt{CountVectorizer}. We kept terms that appear in at least 5 articles (\texttt{min\_df}$=5$) and in no more than 50\% of articles (\texttt{max\_df}$=0.5$). The resulting document-term matrix had 21,528 terms and is 98.97\% sparse (221 distinct terms per article on average), falling within the expected sparsity range for news articles.

Using scikit-learn's \texttt{LatentDirichletAllocation}, we fit and scored models with $K\in\{10,20,30,40,50,60,70\}$ topics using seed 42 (Table~\ref{tab:lda-grid}). Because perplexity on the training articles generally falls as $K$ increases, selecting $K$ involved weighing $c_v$, perplexity, and diversity, as well as manually reviewing selected topics.

\todo{[Laura] Finish the stability-check paragraph (seeds 7 and 21).}

\begin{table}[H]
\caption{LDA topic-count study (seed 42; 7,332 articles). $c_v$ is the mean over topics, using the top ten unigrams per topic and the shared reference corpus. Perplexity is computed on the training articles (lower is better). Diversity is the share of unique terms among all topics' top ten. Fit is wall-clock training time. Bold: best value in each column. Preliminary.}
\label{tab:lda-grid}
\centering
\footnotesize
\setlength{\tabcolsep}{3pt}
\begin{tabular}{rrrrr}
\toprule
\textbf{Topics ($K$)} & \textbf{$c_v$} & \textbf{Perplexity} & \textbf{Diversity} & \textbf{Fit (s)} \\
\midrule
10 & \textbf{0.651} & 3518 & \textbf{0.830} & 29 \\
20 & 0.627 & 3354 & 0.765 & 37 \\
30 & 0.622 & 3323 & 0.720 & 46 \\
40 & 0.647 & 3309 & 0.698 & 49 \\
50 & 0.641 & \textbf{3288} & 0.698 & 55 \\
60 & 0.641 & 3314 & 0.698 & 64 \\
70 & 0.643 & 3355 & 0.689 & 68 \\
\bottomrule
\end{tabular}
\end{table}



\todo{[Laura] Add the LDA formulation: $\theta_d \sim \mathrm{Dir}(\alpha)$, $\phi_k \sim \mathrm{Dir}(\eta)$, $z_{dn} \sim \mathrm{Mult}(\theta_d)$, $w_{dn} \sim \mathrm{Mult}(\phi_{z_{dn}})$; the dominant topic is $\arg\max_k \theta_{dk}$. Name the implementation (gensim \texttt{LdaModel} with \texttt{alpha='auto', eta='auto'} to match ``learned $\alpha,\eta$''), passes and iterations, and seeds.}
\todo{[Laura] Briefly justify LDA as the baseline: interpretable mixtures, a standard comparison in the topic-modeling literature, and low cost. Mark each element as implemented or planned.}

\subsection*{Modern Pipeline (BERTopic)}
The modern pipeline follows five stages \cite{alammar2024, kumari2026}: (1) data preprocessing, (2) document embedding, (3) dimensionality reduction, (4) document clustering, and (5) topic representation and interpretation. We embed articles with GTE-small~\cite{li2023gte} to capture contextual similarity, reduce dimensionality with UMAP~\cite{mcinnes2018umap}, cluster with HDBSCAN~\cite{campello2013hdbscan}, and extract topic keywords with class-based TF-IDF. Embedding input retains sentence context after basic cleaning, and long articles are embedded in chunks that are combined into a single article representation. Because HDBSCAN allows articles to remain unassigned, we report the unassigned proportion explicitly. We examine a limited set of UMAP neighborhood sizes and HDBSCAN minimum cluster sizes.

\textbf{Document embedding.} GTE-small accepts at most 512 tokens. We tokenize each article's full text with the model's tokenizer and split it into consecutive, non-overlapping chunks of at most 480 tokens. Each chunk is wrapped in the model's start and end tokens and embedded separately, giving chunk vectors $\mathbf{v}_{d1},\dots,\mathbf{v}_{dJ}\in\mathbb{R}^{384}$ with token counts $n_{d1},\dots,n_{dJ}$. The article embedding is the token-weighted mean of its chunk vectors, scaled to unit length:
\begin{equation}
\bar{\mathbf{v}}_d = \frac{\sum_{j=1}^{J} n_{dj}\,\mathbf{v}_{dj}}{\sum_{j=1}^{J} n_{dj}}, \qquad
\mathbf{e}_d = \frac{\bar{\mathbf{v}}_d}{\lVert\bar{\mathbf{v}}_d\rVert_2}.
\label{eq:pooling}
\end{equation}
The model revision is pinned, and embeddings are cached so that every configuration uses identical vectors.

\textbf{Reduction and clustering.} UMAP reduces the embeddings to 5 dimensions (cosine metric, \texttt{min\_dist}$=0$), because density-based clustering degrades in high dimensions. HDBSCAN then clusters the reduced vectors (Euclidean metric, \texttt{min\_samples}$=10$, excess-of-mass cluster selection). HDBSCAN infers the number of topics from the data rather than requiring it in advance, and places articles that fit no dense cluster in an outlier class ($-1$) instead of forcing them into a topic.

\textbf{Topic representation.} To describe each topic, BERTopic concatenates the articles assigned to a cluster $c$ into one class document and scores each term $t$ with class-based TF-IDF~\cite{grootendorst2022bertopic}:
\begin{equation}
W_{t,c} = \mathrm{tf}_{t,c}\cdot\log\!\left(1+\frac{A}{f_t}\right),
\label{eq:ctfidf}
\end{equation}
where $\mathrm{tf}_{t,c}$ is the frequency of $t$ in class $c$ (normalized by the class's total term count), $f_t$ is the frequency of $t$ across all classes, and $A$ is the average number of terms per class. A topic's keywords are its ten highest-weighted terms. Counts use unigrams, English stopwords removed, from a fixed vocabulary learned once from the individual articles with \texttt{min\_df}$=5$: a term must appear in at least five articles. Learning the vocabulary before concatenation keeps the threshold defined over articles rather than over class documents, and the same vocabulary is reused across configurations. Keyword tokenization matches the reference corpus used for coherence, so every topic term can be scored.

\subsection*{Topic Labeling and Publication Comparison}
For each pipeline, we label topics using top keywords and representative articles, then produce heat maps of topic proportions within each publication. These comparisons are reported alongside publication sample sizes and the modern model's unassigned proportion. Differences are interpreted descriptively, taking into account unequal publication and date coverage.

\todo{[All] Add an ``Implementation environment'' paragraph (PR1 feedback asks for computing resources): Python 3.13 with the \texttt{uv} lockfile; bertopic 0.17.4, sentence-transformers 6.1.0, gensim 4.4.0, umap-learn 0.5.12, hdbscan 0.8.44, scikit-learn 1.9.1; Apple-silicon macOS (MPS); repository URL.}
\todo{[All] List planned comparisons and ablations, e.g., c-TF-IDF vs.\ alternative keyword representations (\texttt{representation\_comparison.csv}), and LDA with vs.\ without custom stopwords.}

\section*{Experimental Design}
\textbf{Purpose:} Describe how we will develop, evaluate and compare the LDA and BERTopic modeling pipelines.

\textbf{Planned content:} 
\begin{itemize}
    \item Model inputs, metadata and hyperparameters \wip
    \item LDA model development and selection plan \wip
    \item BERTopic model development and selection plan \wip
    \item Independent human evaluation plan \wip
    \item Pipeline comparison \planned
    \item Definition of success \planned
\end{itemize}

\textbf{Working draft:}

\begin{enumerate}
    \item \textbf{Fit both pipelines.} We fit LDA and BERTopic on the same 7,332 articles as described in the Methodology section, documenting preprocessing choices, hyperparameters, and random seeds. For each candidate configuration we record topic counts, topic sizes, and, for BERTopic, the proportion of unassigned articles.
    \item \textbf{Select configurations and label topics.} We considered $c_v$ coherence~\cite{roder2015coherence}, topic sizes, keywords, and representative articles when selecting one configuration per pipeline. Configurations and labels will be finalized before independent human evaluation. Final evaluation articles and duplicate copies will be excluded from development review and topic naming.
    \item \textbf{Independent human evaluation.} Two team members will independently review the same 100 randomly selected articles using a fixed seed. They will rate how well each assigned topic describes the article on a 1--5 scale. We will summarize ratings and reviewer agreement, recording ambiguous and unassigned cases separately.
    \item \textbf{Compare topics across publications.} For each pipeline, we will calculate the proportion of each publication's articles assigned to each topic. Heat maps and representative articles will help identify shared themes and differences in thematic emphasis. We will report publication sample sizes, include BERTopic's unassigned category, and consider unequal date coverage when interpreting differences.
    \item \textbf{Compare pipeline outputs.} We will compare coherence, keyword diversity, human ratings, computational cost, and topic overlap. We will also examine whether both methods support similar conclusions about publication differences, accounting for differences in topic granularity.
    \item \textbf{Assess robustness.} We will repeat selected configurations with different random seeds and examine sensitivity to hyperparameter selection. Where article counts permit, we will restrict publication comparisons to shared periods to assess the influence of unequal date coverage. We will assess whether the main themes and publication patterns persist.
    \item \textbf{Evaluate success and document conclusions.} Success will mean identifying article-supported themes, characterizing their prevalence across publications, and explaining how these findings depend on the modeling method. Findings and conclusions will be expanded once experimentation and analysis are complete.
\end{enumerate}

\todo{[All] Data splitting (assignment item 6): state that there is no train/test split because the task is unsupervised description of a fixed corpus; both models are fit on all 7,332 articles. Leakage is controlled by fixing the human-evaluation sample before labeling.}
\todo{[All] PR1 feedback: define the coherence statistic. $c_v$ (R\"oder et al.) uses the top 10 unigrams per topic, a boolean sliding window of 110 tokens, NPMI with indirect cosine confirmation, and the shared tokenized corpus as reference (gensim \texttt{CoherenceModel}). Report the per-topic values and their mean; BERTopic's $-1$ class is excluded.}
\todo{[All] PR1 feedback: define topic alignment across methods. Jaccard similarity of top-10 term sets, $J(A,B)=|A\cap B|/|A\cup B|$, for every LDA--BERTopic topic pair, shown as a heat map.}
\todo{[All] PR1 feedback: define reviewer agreement. Exact agreement and quadratic-weighted Cohen's $\kappa$ on the 1--5 ratings, per model, over the same 100 articles (seed 42).}
\todo{[All] Hyperparameter selection procedure: BERTopic shortlists the top 3 runs by $c_v$, then picks by mean human rating (semantic consistency, label specificity) of the top 5 topics. LDA: select $K$ by $c_v$ plus inspection of topics.}
\todo{[All] Seeds and uncertainty: grid at seed 42; selected configuration refit with seeds 7 and 21 (ARI for stability). Add bootstrap confidence intervals for mean $c_v$ and for $\kappa$, or state why they are omitted.}
\todo{[All] Compute resources: hardware, cold vs.\ cached embedding time, and per-stage times from the run manifests.}

\section*{Results}
\textbf{Purpose:} Present the themes, publication patterns, and pipeline comparison supported by the completed analysis.

\textbf{Planned content:} 
\begin{itemize}
    \item Topic keywords, labels, sizes and representative articles \planned
    \item Topic proportions across publications, one heat map per pipeline, alongside publication sample sizes \planned
    \item Quantitative and human evaluations, with a comparison table of metrics, human ratings, and computational cost \planned
    \item Agreement and differences between pipelines in themes and publication patterns \planned
    \item Robustness and uncertainty checks \planned
\end{itemize}


\section*{Discussion}
\textbf{Purpose:} Interpret the identified themes and publication differences in relation to the research questions, modeling choices, and prior work on AV news coverage.

\textbf{Planned content:} 
\begin{itemize}
    \item Interpretation of shared themes and differences in thematic emphasis across publications \planned
    \item Influence of pipeline choice on substantive findings \planned
    \item Relationship to prior AV media research \planned
    \item Implications of evaluation results \planned
    \item Discussion of robustness and uncertainty \planned 
    \item Potential extensions and future work \planned
\end{itemize}

\section*{Limitations and Responsible Use Considerations}
\textbf{Purpose:} Acknowledge issues and limitations of our corpus, methods, and evaluation.

\textbf{Planned content:} 
\begin{itemize}
    \item Uneven publications \planned
    \item Imperfect filtering and repeated stories \planned
    \item Influence of major news events \planned
    \item Model assumptions and sensitivity \planned
    \item Limitations of human evaluation \planned
\end{itemize}

\textbf{Working draft:}
The corpus may not represent all AV news coverage. Keyword filtering may omit relevant articles or retain irrelevant ones, while syndicated stories may overrepresent particular themes. Unequal publication sizes and date coverage limit comparisons, especially for publications with few articles. Major news events may also produce temporary concentrations of coverage that do not reflect sustained thematic emphasis.

Both pipelines depend on preprocessing and modeling choices. LDA's dominant-topic summaries simplify articles containing multiple themes, while BERTopic's embeddings, chunk pooling, and clustering settings may affect topic boundaries and leave articles unassigned. Robustness checks will assess these effects. Human evaluation by two team members provides a limited, subjective assessment.

\section*{Conclusion}
\textbf{Purpose:} Summarize the findings about AV news themes and publication differences, explain the contribution of comparing modeling pipelines, and identify remaining work.

\textbf{Planned content:} 
\begin{itemize}
    \item Main themes and differences across publications \planned
    \item Influence of pipeline choice on findings \planned
    \item Remaining analysis and future work \planned
\end{itemize}

\begin{center}\large\textbf{Response to Feedback}\end{center}
\todo{[All] Required section (assignment item 9). Briefly explain how each PR1 comment was addressed and where in the paper:}
\begin{itemize}
    \item \todo{Coherence statistic, topic alignment, reviewer agreement, chunk combination: now defined (Experimental Design; Modern Pipeline).}
    \item \todo{Duplicate handling reconciled: completed; 5 duplicate bodies dropped, earliest copy kept (Data).}
    \item \todo{Primary citations for BERTopic and GTE, plus computing resources (Related Work; Implementation environment).}
    \item \todo{Kang et al.\ (2025) discussed and contribution emphasized (Related Work): done.}
    \item \todo{Major changes: corpus is now 7,332 articles; LDA priors are learned with gensim; actual BERTopic grid differs from the one proposed.}
\end{itemize}

\begin{center}\large\textbf{Remaining Work}\end{center}
\todo{[All] Required section (assignment item 10). Give a specific plan to Progress Report 3, covering:}
\begin{itemize}
    \item \todo{Experiments: LDA grid and seed runs (Laura); BERTopic topic labels (Mary Kate); joint Jaccard and publication heat maps (all); two-reviewer human evaluation (two members).}
    \item \todo{Technical problems: corpus hash changed after duplicate removal, so cached embeddings and runs must be regenerated; finalize custom stopwords.}
    \item \todo{Literature and analysis: primary citations listed above; date-restricted robustness check.}
    \item \todo{Tables and figures still to produce: LDA results table, Jaccard heat map, publication-proportion heat maps for both models, human-evaluation table.}
    \item \todo{Responsibilities: Laura, classical pipeline and evaluation; Mary Kate, modern pipeline and evaluation; Lauren, theory writing and R visualizations.}
    \item \todo{Timeline: Nov 18, seed stability and joint quantitative and human evaluation; Dec 2, report tables, figures, methods and results, presentation. Note any scope changes.}
\end{itemize}

\newpage
\section*{References}
\printbibliography[heading=none]

\end{document}